En physique, 
on apprend à modéliser le réel pour en comprendre les règles. 

À force d'aborder des problèmes complexes en TD, 
j'ai réalisé que ce qui m'intéressait le plus n'était pas seulement d'observer ces phénomènes, 
mais d'utiliser ces outils pour \textbf{concevoir des solutions opérationnelles}. 

C'est ce constat qui oriente aujourd'hui mon projet : 
faire le pont entre mes acquis académiques que je continue d'accumuler
et les \textbf{exigences du métier d'ingénieur}. 

Pour y parvenir, 
je ne conçois pas ce parcours que j'ai choisi comme une simple fin en soi, 
mais comme le moyen de \textbf{construire mon employabilité}. 

À travers ce document, 
je vous propose de suivre ma trajectoire : 
\begin{itemize}
    \item \textbf{Pourquoi ce virage :} de la passion scientifique au métier d'ingénieur ;
    \item \textbf{Pourquoi ce choix :} une vision plus argumentée basée sur les besoins réels du marché du travail;
    \item \textbf{Ma feuille de route :} transformer mon savoir scientifique en levier d'employabilité ;
    \item \textbf{Cap sur l'avenir :} ma vision d'ingénieur en action.
\end{itemize}